\documentclass[10pt,a4paper]{amsart}
\usepackage[margin=19mm]{geometry}
\usepackage{amsmath,amssymb,mathtools}
\usepackage[scr=rsfs,cal=euler]{mathalfa}
\usepackage{xfrac}
\usepackage[unicode,hidelinks,bookmarks=false]{hyperref}
\newcommand{\ie}{i.e.\ }
\newcommand{\cf}{cf.\ }
\newcommand{\resp}{respectively\ }
\newcommand{\Subsection}{\subsection}
\providecommand{\nobreakdash}{\nobreak\hbox{-}}
\setlength{\emergencystretch}{2em}
\multlinegap0pt
\usepackage{bm}
\usepackage{bbm}
\usepackage{dsfont}
\usepackage{xspace}
\usepackage{cancel}
\usepackage{mathtools}
\usepackage{amsthm}
\makeatletter
\@ifundefined{theorem}{\newtheorem{theorem}{Theorem}}{}
\@ifundefined{lemma}{\newtheorem{lemma}[theorem]{Lemma}}{}
\makeatother
\def\C{\mathbb{C}}
\def\H{\mathbb{H}}
\def\g{\gamma}
\title[Warped product metrics on hyperbolic and complex hyperbolic ]{Warped product metrics on hyperbolic and complex hyperbolic manifolds}
\author{Barry Minemyer}
\date{Source: arXiv:2307.15553v1 (CC BY 4.0)}
\begin{document}
\maketitle

\noindent\emph{Source and licence.} Warped product metrics on hyperbolic and complex hyperbolic manifolds. Version arXiv:2307.15553v1 (CC BY 4.0), distributed under the Creative Commons Attribution 4.0 International licence, as recorded in the arXiv licence metadata for this exact version and shown on the abstract page https://arxiv.org/abs/2307.15553v1. Full rights notice: licenses/arxiv-2307.15553-CC-BY-4.0.txt.\par\smallskip

\noindent\textit{Editorial scope.} Counted unit: the Lemma whose source label is \texttt{None} in the article (source0.tex lines 1055--1057, source label \texttt{None}), delivered together with the article's own complete original proof of it (source0.tex lines 1059--1094, one \texttt{proof} environment of 36 lines). No mathematics is written, repaired or re-derived: statement and proof are the article's own text line for line. Required context retained uncounted: eqn:/g (lines 797--799); eqn:orthonormal basis 2 (lines 814--817); thm:curvature tensor 2 (lines 1007--1024). Every definition, display and label of those lines is retained unchanged. Omitted from this package and not redistributed: 1--796, 800--813, 818--1006, 1025--1054, 1058--1058, 1095--1524. Nothing inside a retained excerpt is omitted or altered: every retained body is delivered exactly as the article's lines are written, so no omission receipt of any kind accompanies this package. Citations: the retained excerpts carry no citation command at all, so no bibliography entry of the article is transcribed and no reference is redistributed. Reference adaptation: none was needed and none was made; every reference inside the delivered unit resolves inside it, so no unresolved reference or citation is produced. Printed slips kept literal: no printed word, symbol, hypothesis or number of the counted statement or of its proof was altered. Layout and class: the article's own class is \texttt{amsart} with packages including latexsym, amsmath, amsfonts, amssymb, mathrsfs, fancyhdr, tikz, tikz-cd, verbatim, multicol; the wrapper is the lane's pinned \texttt{amsart} editorial wrapper at 10pt on A4, which restates the theorem-like environments the retained text needs and adds the article's own macro definitions that the retained text needs (\texttt{\textbackslash C}, \texttt{\textbackslash H}, \texttt{\textbackslash g}), with extra wrapper packages (bm, bbm, dsfont, xspace, cancel, mathtools). The delivered numbering is the unit's own. Assets: no retained span contains a figure environment, a graphics-inclusion command, a PDF-page inclusion, a table or a TikZ picture; no image file is copied into this package, the package declares no asset dependency and no image file is redistributed with it. Licence: version arXiv:2307.15553v1 of this article is distributed under the Creative Commons Attribution 4.0 International licence, as recorded in the arXiv licence metadata for this exact version and shown on the abstract page https://arxiv.org/abs/2307.15553v1. A full rights notice is delivered at licenses/arxiv-2307.15553-CC-BY-4.0.txt. Characters: every retained excerpt is pure ASCII, so the delivered engine, XeLaTeX with its default Latin Modern text font, reports no missing character.
	\begin{equation}\label{eqn:/g}
	{\bf \g} := {\bf \g_r} + dr^2.
	\end{equation}

	\begin{align}
	&Y_1 = \frac{1}{h} X_1	\hskip 25pt	Y_2 = \frac{1}{h} X_2	\hskip 25pt	Y_3 = \frac{1}{h} X_3		\hskip 25pt	Y_4 = \frac{1}{h} X_4  \hskip 25pt  Y_5 = \frac{1}{v} X_5  \label{eqn:orthonormal basis 2}	  \\
	&Y_6 = \frac{1}{v} X_6		\hskip 25pt	Y_7 = \frac{1}{v} X_7		\hskip 25pt	Y_8 = \frac{1}{v} X_8	\hskip 25pt	Y_9 = \frac{1}{\frac{1}{2} v_r} X_9  \hskip 25pt  Y_{10} = X_{10}.  \nonumber
	\end{align}

\begin{theorem}\label{thm:curvature tensor 2}
In terms of the basis given in equation \eqref{eqn:orthonormal basis 2}, the only independent nonzero components of the $(4,0)$ curvature tensor $R_{\g}$ are given by the following formulas, where $i \in \{ 1, 2, 3, 4 \}$, $k \in \{ 5, 6, 7, 8 \}$, $( i, j) \in \{ (1, 2), ( 3, 4) \}$, and $(k, l) \in \{ (5, 6), ( 7, 8) \}$.

	\begin{align*}
	&R^{\g}_{1212} = R^{\g}_{3434} = - \left( \frac{h'}{h} \right)^2 - \frac{4}{h^2} - \frac{3 v_r^2}{4h^4}  \hskip 20pt R^{\g}_{5656} = R^{\g}_{7878} = - \left( \frac{v'}{v} \right)^2 + \frac{4}{v^2} - \frac{3v_r^2}{4 v^4}  \\
	&R^{\g}_{i9i9} = - \frac{h' v_r'}{h v_r} + \frac{v_r^2}{4 h^4}   \hskip 40pt	R^{\g}_{k9k9} = - \frac{v' v_r'}{v v_r} + \frac{v_r^2}{4 v^4}	\hskip 40pt	R^{\g}_{ikik} = - \frac{h' v'}{hv}   \\
	&R^{\g}_{1313} = R^{\g}_{1414} = R^{\g}_{2323} = R^{\g}_{2424} = - \left( \frac{h'}{h} \right)^2 - \frac{1}{h^2}  \hskip 40pt	\text{(continued on next page)}
	\end{align*}
	\begin{align*}
	&R^{\g}_{5757} = R^{\g}_{5858} = R^{\g}_{6767} = R^{\g}_{6868} = - \left( \frac{v'}{v} \right)^2 + \frac{1}{v^2}  \\
	&R^{\g}_{i,10,i,10} = - \frac{h''}{h}  \hskip 60pt	R^{\g}_{k,10,k,10} = - \frac{v''}{v} 	\hskip 60pt R^{\g}_{9,10,9,10} = - \frac{v_r''}{v_r}   \\
	&R^{\g}_{1234} = 2R^{\g}_{1324} = -2R^{\g}_{1423} = - \frac{2}{h^2} - \frac{v_r^2}{2h^4}  \\
	&R^{\g}_{5678} = 2R^{\g}_{5768} = -2R^{\g}_{5867} = \frac{2}{v^2} - \frac{v_r^2}{2v^4}  \\
	&R^{\g}_{ijkl} = 2R^{\g}_{ikjl} = -2R^{\g}_{iljk} = - \frac{v_r^2}{2 h^2 v^2}   \\
	&R^{\g}_{i,j,9, 10} = 2R^{\g}_{i,9,j,10} = -2R^{\g}_{i,10,j,9} = - \frac{v_r}{h^2} \left( \ln \frac{v_r}{h} \right)'   \\
	&R^{\g}_{k,l,9,10} = 2R^{\g}_{k,9,l,10} = -2R^{\g}_{k,10,l,9} = - \frac{v_r}{v^2} \left( \ln \frac{v_r}{v} \right)' .
	\end{align*}
\end{theorem}

\begin{lemma}
There do not exist functions $h, v, $ and $v_r$ that, when inserted into equation \eqref{eqn:/g}, yield a complete finite volume Riemannian metric on $\C \H^n \setminus \C \H^{n-2}$ with nonpositive sectional curvature.
\end{lemma}

\begin{proof}
Since we are removing a copy of $\C \H^{n-2}$ from $\C \H^n$, for the metric to be complete we need to define $h, v, $ and $v_r$ on $(- \infty, \infty)$.  
These functions need to be positive for the metric to be Riemannian, and they need to be non-decreasing for there to be any chance of nonpositive curvature.  
For any hope of finite volume, we need all of the following limits to hold:
	\begin{equation*}
	\lim_{r \to - \infty} h, h', v, v', v_r, v_r' = 0.
	\end{equation*}
	
Now, from the formula for the $R_{5656}$ term in Theorem \ref{thm:curvature tensor 2} we must have that
	\begin{align}
	\frac{4-(v')^2}{v^2} - \frac{3v_r^2}{4v^4} &\leq 0  \label{eqn:neg curv 1}  \\
	\iff	\hskip 30pt	\frac{16 - 4(v')^2}{3} &\leq \left( \frac{v_r}{v} \right)^2.  \nonumber
	\end{align}
In particular, we see that a necessary requirement for nonpositive curvature is that
	\begin{equation}\label{eqn:contradiction 1}
	\lim_{r \to - \infty} \frac{v_r}{v} > 1.
	\end{equation}
	
From the formula for the $R_{k9k9}$ term in Theorem \ref{thm:curvature tensor 2} we must have that
	\begin{align}
	- \frac{v' v_r'}{v v_r} + \frac{v_r^2}{4v^4} &\leq 0  \nonumber  \\
	\Longrightarrow	\hskip 30pt	\frac{3v_r^2}{4v^4} &\leq \frac{3 v' v_r'}{v v_r}.  \label{eqn:neg curv 2}
	\end{align}
Comparing equations \eqref{eqn:neg curv 1} and \eqref{eqn:neg curv 2}, we see that
	\begin{equation*}
	\frac{4 - (v')^2}{v^2} \leq \frac{3 v' v_r'}{v v_r}	\hskip 20pt	\Longrightarrow		\hskip 20pt	4 - (v')^2 \leq (3v'v_r') \cdot \frac{v}{v_r}
	\end{equation*}
is also a necessary requirement for nonpositive curvature.  
But as $r \to - \infty$, we know that $4 - (v')^2 \to 4$ and $3 v' v_r' \to 0$.  
Thus, we must have that
	\begin{equation}\label{eqn:contradiction 2}
	\lim_{r \to - \infty} \frac{v}{v_r} = \infty	\hskip 20pt	\Longrightarrow		\hskip 20pt	\lim_{r \to - \infty} \frac{v_r}{v} = 0.
	\end{equation}
Equations \eqref{eqn:contradiction 1} and \eqref{eqn:contradiction 2} provide a contradiction, proving the Lemma.
	
\end{proof}

\end{document}
